\documentclass[11pt,a4paper]{article}
\usepackage[margin=2.2cm]{geometry}
\usepackage[utf8]{inputenc}
\usepackage[T1]{fontenc}
\usepackage{amsmath,amssymb,amsfonts,bm}
\usepackage{braket}
\usepackage{enumitem}
\usepackage{xcolor}
\usepackage{tikz}
\usetikzlibrary{quantikz2}

\newcommand{\pts}[1]{\textbf{(#1)}}
\newcommand{\tr}{\operatorname{tr}}
\newcommand{\id}{\mathbb{I}}
\newcommand{\sol}{\color{blue!60!black}}

\setlist[enumerate,1]{label=\alph*)}
\setlength{\parindent}{0pt}
\setlength{\parskip}{4pt}

\title{Quantum Technologies -- First partial exam\\[4pt] \large Practice exam A}
\date{}

\begin{document}
\maketitle
\vspace{-1.5cm}
\begin{center}
Duration: 2 hours. Total: 10 points.
\end{center}

%=====================================================================
\section{Open question (1pt)}

A quantum computer can evaluate a function $f:\{0,1\}^n\to\{0,1\}$ on all $2^n$ inputs with a
single call to the oracle $U_f$ (\emph{quantum parallelism}). Explain why quantum parallelism by itself
does not provide a computational advantage, and which additional ingredient is needed to extract
useful information. Illustrate your answer with TWO quantum algorithms seen in class, indicating
in each case what the algorithm learns about $f$ and the speed-up with respect to the best
classical strategy. (Maximum 200 words, $\sim$1 page)

%=====================================================================
\section{Short questions (3pts)}

\begin{enumerate}
\item Consider the following circuit acting on three qubits. Can it be used as the oracle of the
Grover algorithm? If so, indicate (i) which is[are] the marked item[s] and (ii) the number of
Grover iterations that maximizes the probability of finding a marked item, together with that
probability \pts{1 p}.
\begin{center}
\begin{quantikz}[row sep=0.25cm, column sep=0.35cm]
\lstick{$\ket{q_0}$} & \gate{X} & \ctrl{2} & \gate{X} & \\
\lstick{$\ket{q_1}$} & & & & \\
\lstick{$\ket{q_2}$} & & \gate{Z} & &
\end{quantikz}
\end{center}

\item We use the quantum phase estimation algorithm to find the phases $\varphi$ associated to the
eigenstates $\ket{u}=\ket{0}$ and $\ket{u}=\ket{1}$ of the unitary operator
$U = R_z(\pi/2) = \begin{pmatrix} e^{-i\pi/4} & 0\\ 0 & e^{i\pi/4}\end{pmatrix}$.
If we use $n=4$ qubits in the first register, what is the expected measurement $\vec{x}$ for each
qubit in the register, for each of the two eigenstates? \pts{1 p}

\item Solving the order-finding problem with a quantum computer for $x=2$ and $M=21$, we measure
the value $\vec{y}_0=\{001011\}$ in the first register, formed by $n=6$ qubits. With this
information, use the continued fraction expansion to find the order of $x$ modulo $M$. Would the
measurement $\vec{y}_0^{\,\prime}=\{101000\}$ also be compatible with this problem? \pts{1 p}
\end{enumerate}

%=====================================================================
\section*{Problems}

\section{An imperfect state preparation (3pts)}

A device is designed to prepare a qubit in the state $\ket{0}$. Due to experimental imperfections,
it prepares $\ket{0}$ with probability $p$ and the state $\ket{+}=(\ket{0}+\ket{1})/\sqrt{2}$ with
probability $1-p$, with $0\le p\le 1$.

\begin{enumerate}
\item Write the density matrix $\rho$ describing the output of the device in the computational
basis, and check that it satisfies the properties of a density matrix \pts{0.5 p}.
\item Find its Bloch vector $\bm{r}$ and locate it in the Bloch sphere. For which values of $p$ is
$\rho$ a pure state? \pts{0.5 p}
\item Calculate the purity $\gamma(\rho)$ and the eigenvalues of $\rho$ as a function of $p$.
Evaluate the von Neumann entropy for $p=1/2$ \pts{0.5 p}.
\item The observable $\sigma_x$ is measured on the output of the device. Which are the possible
outcomes, their probabilities and the state after the measurement for each outcome? What is the
average value $\langle\sigma_x\rangle$? \pts{0.5 p}
\item Calculate the fidelity $F(\rho,\ket{0}\!\bra{0})$ between the prepared state and the target
state \pts{0.5 p}.
\item Calculate the trace distance $d(\rho,\ket{0}\!\bra{0})$. Is $\ket{0}$ the pure state with the
largest fidelity with $\rho$? If not, which pure state is it, and what is that fidelity? \pts{0.5 p}
\end{enumerate}

\section{Controlled rotation (3pts)}

Consider the following two-qubit circuit, where
$R_y(\theta)=\begin{pmatrix}\cos\frac{\theta}{2} & -\sin\frac{\theta}{2}\\ \sin\frac{\theta}{2} & \cos\frac{\theta}{2}\end{pmatrix}$
and $0\le\theta<2\pi$.
\begin{center}
\begin{quantikz}[row sep=0.4cm, column sep=0.4cm]
\lstick{$\ket{q_0}$} & \gate{H} & \ctrl{1} & \\
\lstick{$\ket{q_1}$} & & \gate{R_y(\theta)} &
\end{quantikz}
\end{center}

\begin{enumerate}
\item Write the matrix of the circuit in the computational basis $\{\ket{00},\ket{01},\ket{10},\ket{11}\}$ \pts{0.5 p}.
\item Consider the input state $\ket{00}$. Write the output state $\ket{\psi}$. For which values of
$\theta$ is the output a product state? For which value is it a Bell state? \pts{0.5 p}
\item Calculate the reduced density matrices $\rho_A$ (first qubit) and $\rho_B$ (second qubit),
their Bloch vectors and their purities. For which values of $\theta$ is $\ket{\psi}$ entangled?
\pts{0.5 p}
\item For $\theta=2\pi/3$, find the Schmidt decomposition of $\ket{\psi}$ and its entanglement
entropy \pts{1 p}.
\item For $\theta=2\pi/3$, we measure $\sigma_z$ on the second qubit. Calculate the probability of
each outcome and the state of the first qubit after the measurement. Is the two-qubit state still
entangled after the measurement? \pts{0.5 p}
\end{enumerate}

\textit{Useful formulas:} $1+\cos x = 2\cos^2\frac{x}{2}$, $\;1-\cos x = 2\sin^2\frac{x}{2}$, $\;2\sin x\cos x=\sin 2x$.

%=====================================================================
\newpage
\begin{center}
\Large\textbf{Solutions}
\end{center}
\sol

\section*{1\quad Open question}

Using $n$ qubits in the superposition $\ket{s}=H^{\otimes n}\ket{0}^{\otimes n}$, one call to
$U_f:\ket{x}\ket{y}\mapsto\ket{x}\ket{y\oplus f(x)}$ produces
$\frac{1}{\sqrt N}\sum_x\ket{x}\ket{f(x)}$, i.e.\ the function is evaluated on all
$N=2^n$ inputs at once. However, the values $f(x)$ are not accessible: a measurement collapses the
state to a single random pair $(x,f(x))$, which is no better than one classical evaluation.

The additional ingredient is \textbf{interference}. Preparing the target in $\ket{-}$, the oracle
writes $f$ into the phases, $\ket{x}\ket{-}\mapsto(-1)^{f(x)}\ket{x}\ket{-}$ (phase kickback). A further
unitary (Hadamards, QFT, diffuser) makes the amplitudes interfere so that the answer to a
\emph{global} question about $f$ is concentrated on a few computational states.
\begin{itemize}
\item \textbf{Deutsch--Jozsa}: decides whether $f$ is constant or balanced. After
$H^{\otimes n}$, the probability of measuring $0\dots0$ is $|\frac1N\sum_x(-1)^{f(x)}|^2$, which is 1
for a constant function and 0 for a balanced one. One query, compared to $2^{n-1}+1$ classical
queries in the worst case (exponential speed-up).
\item \textbf{Bernstein--Vazirani}: for $f(x)=\vec s\cdot\vec x$, the state
$\frac1{\sqrt N}\sum_x(-1)^{\vec s\cdot\vec x}\ket{x}$ is mapped by $H^{\otimes n}$ to $\ket{s}$. The
hidden string is found with one query instead of $n$.
\item (Also valid: \textbf{Grover}, where the oracle marks items with a phase and the diffuser
amplifies their amplitude, giving $\mathcal{O}(\sqrt N)$ queries instead of $\mathcal{O}(N)$.)
\end{itemize}

\section*{2\quad Short questions}

\begin{enumerate}
\item Yes. The gate $X$--$CZ$--$X$ only adds a sign $-1$ when the control (after the first $X$) and
the target are both in $\ket1$, that is when $q_0=0$ and $q_2=1$, whatever the value of $q_1$. It is
diagonal with entries $\pm1$, so it is a valid oracle $U_f\ket{x}=(-1)^{f(x)}\ket{x}$, with
\[
\text{marked items: } \ket{w}\in\{\ket{001},\ket{011}\}\quad(M=2,\ N=2^3=8).
\]
With $M$ marked items, $\sin\theta_0=\sqrt{M/N}=\sqrt{2/8}=1/2$, so $\theta_0=\pi/6$. After $k$
iterations, $P_k=\sin^2[(2k+1)\theta_0]$:
\[
k_{\max}\simeq\left\lfloor\frac{\pi}{4}\sqrt{\frac{N}{M}}\right\rfloor=\left\lfloor\frac{\pi}{2}\right\rfloor=1,
\qquad P_{k=1}=\sin^2\!\left(3\cdot\frac{\pi}{6}\right)=1 .
\]
A single Grover iteration finds one of the marked items with certainty. (A second iteration
would drop the probability back to $\sin^2(5\pi/6)=1/4$.)

\item $U$ is diagonal, so $\ket0$ and $\ket1$ are eigenstates. We write each eigenvalue as
$e^{2\pi i\varphi}$ with $\varphi\in[0,1)$:
\[
U\ket1=e^{i\pi/4}\ket1=e^{2\pi i\frac18}\ket1\ \Rightarrow\ \varphi=\tfrac18=\tfrac{x}{2^4}\ \Rightarrow\ x=2\equiv\{0010\}.
\]
\[
U\ket0=e^{-i\pi/4}\ket0=e^{2\pi i\left(1-\frac18\right)}\ket0\ \Rightarrow\ \varphi=\tfrac78=\tfrac{x}{2^4}\ \Rightarrow\ x=14\equiv\{1110\}.
\]
In both cases $\varphi N$ is an integer, so the result is obtained with probability 1. (A negative
phase has to be mapped to $[0,1)$ first.)

\item $\vec y_0=\{001011\}$ corresponds to $y_0=11$, with $N=2^6=64$. We look for $k/r$ such that
$\left|\frac{y_0}{N}-\frac kr\right|<\frac1{2N}=\frac1{128}$. The continued fraction expansion is
\[
\frac{11}{64}=\cfrac{1}{5+\cfrac{1}{1+\cfrac{1}{4+\cfrac12}}}\equiv[0;5,1,4,2],
\qquad\text{convergents: } \frac01,\ \frac15,\ \frac16,\ \frac5{29},\ \frac{11}{64}.
\]
$\left|\frac{11}{64}-\frac15\right|=\frac{9}{320}\approx0.028>\frac1{128}$, while
$\left|\frac{11}{64}-\frac16\right|=\frac{2}{384}\approx0.0052<\frac1{128}\approx0.0078$. Therefore
$r=6$. Check: $2^1,\dots,2^6 \bmod 21 = 2,4,8,16,11,1$, so indeed $2^6\equiv1\bmod 21$ and $r=6$ is the order.

$\vec y_0^{\,\prime}=\{101000\}$ corresponds to $y_0'=40$, $\frac{40}{64}=\frac58=[0;1,1,1,2]$ with
convergents $0,\ 1,\ \frac12,\ \frac23,\ \frac58$. The fractions with denominator $r=6$ closest to
$5/8=0.625$ are $\frac36=0.5$ and $\frac46=0.667$, and
$\left|\frac58-\frac46\right|=\frac1{24}\approx0.042>\frac1{128}$. So $y_0'$ is \textbf{not}
compatible with $r=6$. (The likely outcomes are $y_0\approx 64k/6$: $0,11,21,32,43,53$.)
\end{enumerate}

\section*{3\quad An imperfect state preparation}

\begin{enumerate}
\item The ensemble $\{(p,\ket0),(1-p,\ket+)\}$ gives
\[
\rho=p\ket0\!\bra0+(1-p)\ket+\!\bra+
= p\begin{pmatrix}1&0\\0&0\end{pmatrix}+\frac{1-p}{2}\begin{pmatrix}1&1\\1&1\end{pmatrix}
=\frac12\begin{pmatrix}1+p&1-p\\1-p&1-p\end{pmatrix}.
\]
It is Hermitian (a real symmetric matrix), $\tr\rho=\frac12[(1+p)+(1-p)]=1$, and positive
semi-definite: $\tr\rho=1>0$ and $\det\rho=\frac14[(1+p)(1-p)-(1-p)^2]=\frac{p(1-p)}{2}\ge0$, so both
eigenvalues are non-negative. (It is also a convex combination of projectors, which is positive by construction.)

\item Using $r_\alpha=\tr[\rho\sigma_\alpha]$:
\[
r_x=2\,\mathrm{Re}\,\rho_{01}=1-p,\qquad r_y=-2\,\mathrm{Im}\,\rho_{01}=0,\qquad r_z=\rho_{00}-\rho_{11}=p,
\]
so $\bm r=(1-p,\,0,\,p)$, i.e.\ $\rho=\frac12\left[\id+(1-p)\sigma_x+p\,\sigma_z\right]$. This is
the convex combination $p\,\hat{\bm z}+(1-p)\hat{\bm x}$: the segment joining the north pole
($\ket0$) and the point $+\hat{\bm x}$ of the equator ($\ket+$), in the $xz$ plane. Its norm is
$|\bm r|=\sqrt{p^2+(1-p)^2}$, and $|\bm r|^2=1-2p(1-p)=1\iff p(1-p)=0$. So the state is pure only
for $p=0$ ($\rho=\ket+\!\bra+$) or $p=1$ ($\rho=\ket0\!\bra0$). Otherwise it lies inside the Bloch ball.

\item Purity:
\[
\gamma(\rho)=\tr\rho^2=\frac{1+|\bm r|^2}{2}=\frac{1+p^2+(1-p)^2}{2}=1-p(1-p).
\]
Eigenvalues: $\lambda_\pm=\frac12(1\pm|\bm r|)=\frac12\left(1\pm\sqrt{p^2+(1-p)^2}\right)$. One can
check $\lambda_+^2+\lambda_-^2=\frac{1+|\bm r|^2}{2}=\gamma$.

For $p=1/2$: $|\bm r|=1/\sqrt2$, $\lambda_\pm=\frac12\left(1\pm\frac1{\sqrt2}\right)\approx0.854,\ 0.146$, $\gamma=3/4$, and
\[
S(\rho)=-\lambda_+\log_2\lambda_+-\lambda_-\log_2\lambda_-\approx-0.854\log_2 0.854-0.146\log_2 0.146\approx0.60 .
\]

\item $\sigma_x=\ket+\!\bra+-\ket-\!\bra-$, so the possible outcomes are $\pm1$, with projectors
$P_\pm=\ket\pm\!\bra\pm$:
\[
p(+1)=\tr[P_+\rho]=\bra+\rho\ket+=\frac{1+r_x}{2}=1-\frac p2,\qquad
p(-1)=\bra-\rho\ket-=\frac{1-r_x}{2}=\frac p2 .
\]
(Directly: $\bra+\rho\ket+=\frac14[(1+p)+2(1-p)+(1-p)]=1-\frac p2$.) After the measurement the state
is $P_\pm\rho P_\pm/p(\pm1)=\ket\pm\!\bra\pm$ for outcome $\pm1$. The average value is
$\langle\sigma_x\rangle=\tr[\rho\sigma_x]=p(+1)-p(-1)=1-p=r_x$.

\item Since $\ket0\!\bra0$ is pure, $F(\rho,\ket0\!\bra0)=\bra0\rho\ket0=\rho_{00}=\dfrac{1+p}{2}$.
Equivalently, $F=\frac{1+\bm r\cdot\hat{\bm z}}{2}=\frac{1+p}{2}$. It equals 1 only for $p=1$, and it is $1/2$ for
$p=0$, since $|\bra0+\rangle|^2=1/2$.

\item For qubits, the trace distance is half the Euclidean distance between the Bloch vectors.
With $\bm s=(0,0,1)$ for $\ket0$:
\[
d(\rho,\ket0\!\bra0)=\frac12|\bm r-\bm s|=\frac12\sqrt{(1-p)^2+(p-1)^2}=\frac{1-p}{\sqrt2}.
\]
For a pure state $\ket\phi$ with Bloch vector $\hat{\bm n}$, $F=\bra\phi\rho\ket\phi=\frac{1+\bm r\cdot\hat{\bm n}}{2}$.
This is maximal when $\hat{\bm n}\parallel\bm r$, i.e.\ for the eigenvector of $\rho$ with the
largest eigenvalue, giving
\[
F_{\max}=\lambda_+=\tfrac12\left(1+\sqrt{p^2+(1-p)^2}\right).
\]
This is
larger than $\frac{1+p}{2}$ whenever $0\le p<1$, because $\sqrt{p^2+(1-p)^2}>p$. So $\ket0$ is
\textbf{not} the closest pure state. The closest one points along $\bm r$, between $\ket0$ and
$\ket+$. For $p=1/2$, it is the state along $(\hat{\bm x}+\hat{\bm z})/\sqrt2$, with
$F_{\max}\approx0.854>0.75$.
\end{enumerate}

\section*{4\quad Controlled rotation}

Notation: $c\equiv\cos\frac\theta2$, $s\equiv\sin\frac\theta2$.

\begin{enumerate}
\item The circuit is $U=C\text{-}R_y(\theta)\,(H\otimes\id)$, with
$C\text{-}R_y=\begin{pmatrix}\id&0\\0&R_y(\theta)\end{pmatrix}$ and
$H\otimes\id=\frac1{\sqrt2}\begin{pmatrix}\id&\id\\ \id&-\id\end{pmatrix}$ ($2\times2$ blocks). Therefore
\[
U=\frac1{\sqrt2}\begin{pmatrix}\id&\id\\ R_y&-R_y\end{pmatrix}
=\frac1{\sqrt2}\begin{pmatrix}1&0&1&0\\0&1&0&1\\ c&-s&-c&s\\ s&c&-s&-c\end{pmatrix}.
\]
(Equivalently, one can follow the basis states: $\ket{00}\to\ket{+}\ket0\to\frac1{\sqrt2}[\ket{00}+c\ket{10}+s\ket{11}]$, and so on.)

\item The output is the first column of $U$:
\[
\ket\psi=\frac1{\sqrt2}\left[\ket0\ket0+\ket1 R_y(\theta)\ket0\right]
=\frac1{\sqrt2}\left[\ket{00}+c\ket{10}+s\ket{11}\right].
\]
For $\theta=0$ we get $R_y=\id$ and $\ket\psi=\frac1{\sqrt2}(\ket{00}+\ket{10})=\ket+\ket0$, a product
state. For $\theta=\pi$ we get $c=0$, $s=1$, and $\ket\psi=\frac1{\sqrt2}(\ket{00}+\ket{11})=\ket{\beta_{00}}$, a Bell state.

\item We write $\ket\psi=\sum_{ij}\alpha_{ij}\ket{ij}$ with
$\alpha=\frac1{\sqrt2}\begin{pmatrix}1&0\\c&s\end{pmatrix}$. Then $\rho_A=\alpha\alpha^\dagger$ and $\rho_B=\alpha^T\alpha^*$:
\[
\rho_A=\frac12\begin{pmatrix}1&c\\c&c^2+s^2\end{pmatrix}=\frac12\begin{pmatrix}1&c\\c&1\end{pmatrix},
\qquad
\rho_B=\frac12\begin{pmatrix}1+c^2&cs\\cs&s^2\end{pmatrix}.
\]
Bloch vectors ($r_x=2\mathrm{Re}\rho_{01}$, $r_z=\rho_{00}-\rho_{11}$):
\[
\bm r_A=(c,0,0)=\left(\cos\tfrac\theta2,0,0\right),\qquad
\bm r_B=\left(cs,\,0,\,\tfrac{1+c^2-s^2}{2}\right)=\left(\tfrac12\sin\theta,\,0,\,\cos^2\tfrac\theta2\right).
\]
Both have the same norm: $|\bm r_B|^2=c^2s^2+c^4=c^2=|\bm r_A|^2$. So both purities are equal, as
expected for a pure global state:
\[
\gamma(\rho_A)=\gamma(\rho_B)=\frac{1+\cos^2(\theta/2)}{2}.
\]
$\ket\psi$ is entangled $\iff$ $\rho_A$ is mixed $\iff$ $\cos^2(\theta/2)<1$ $\iff$ $\theta\neq0$.
It is maximally entangled ($\gamma=1/2$, $\rho_A=\id/2$) for $\theta=\pi$.

\item For $\theta=2\pi/3$: $c=\frac12$, $s=\frac{\sqrt3}2$, and
$\ket\psi=\frac1{\sqrt2}\left[\ket{00}+\frac12\ket{10}+\frac{\sqrt3}2\ket{11}\right]$.

\emph{Schmidt coefficients and basis of $A$.} $\rho_A=\frac12\begin{pmatrix}1&1/2\\1/2&1\end{pmatrix}=\frac12\left(\id+\frac12\sigma_x\right)$
has eigenvalues $\eta_1^2=\frac34$ (eigenvector $\ket+$) and $\eta_2^2=\frac14$ (eigenvector $\ket-$).
So $\eta_1=\frac{\sqrt3}2$, $\eta_2=\frac12$, with $\ket{1_A}=\ket+$ and $\ket{2_A}=\ket-$.

\emph{Basis of $B$.} We project $\ket\psi$ onto the Schmidt vectors of $A$:
\begin{align*}
\bra{+}_A\ket\psi&=\frac12\left[\ket0+\tfrac12\ket0+\tfrac{\sqrt3}2\ket1\right]=\frac{\sqrt3}2\cdot\underbrace{\frac{\sqrt3\ket0+\ket1}{2}}_{\ket{1_B}},\\
\bra{-}_A\ket\psi&=\frac12\left[\ket0-\tfrac12\ket0-\tfrac{\sqrt3}2\ket1\right]=\frac12\cdot\underbrace{\frac{\ket0-\sqrt3\ket1}{2}}_{\ket{2_B}}.
\end{align*}
$\ket{1_B},\ket{2_B}$ are orthonormal. They are the eigenvectors of
$\rho_B=\frac18\begin{pmatrix}5&\sqrt3\\\sqrt3&3\end{pmatrix}$ with eigenvalues $3/4$ and $1/4$. Therefore
\[
\boxed{\ket\psi=\frac{\sqrt3}{2}\ket{+}\otimes\frac{\sqrt3\ket0+\ket1}{2}+\frac12\ket{-}\otimes\frac{\ket0-\sqrt3\ket1}{2}}
\]
The Schmidt rank is 2, so the state is entangled. The entanglement entropy is
\[
S_{EE}=-\sum_k\eta_k^2\log_2\eta_k^2=-\tfrac34\log_2\tfrac34-\tfrac14\log_2\tfrac14=2-\tfrac34\log_2 3\approx0.81 .
\]
(For general $\theta$: $\eta_1=\cos\frac\theta4$, $\eta_2=\sin\frac\theta4$, $\ket{1_B}=\cos\frac\theta4\ket0+\sin\frac\theta4\ket1$, $\ket{2_B}=\sin\frac\theta4\ket0-\cos\frac\theta4\ket1$.)

\item We group the terms according to the state of the second qubit:
\[
\ket\psi=\frac1{\sqrt2}\left[\left(\ket0+\tfrac12\ket1\right)\ket0+\tfrac{\sqrt3}2\ket1\ket1\right].
\]
With the projectors $P_m=\id\otimes\ket m\!\bra m$:
\[
p(+1)=p(\ket0_B)=\frac12\left(1+\frac14\right)=\frac58,\qquad
\ket{\psi'_A}=\frac{\ket0+\frac12\ket1}{\sqrt{5/4}}=\frac{2\ket0+\ket1}{\sqrt5};
\]
\[
p(-1)=p(\ket1_B)=\frac12\cdot\frac34=\frac38,\qquad \ket{\psi'_A}=\ket1 .
\]
After the measurement the global state is $\ket{\psi'_A}\otimes\ket m$, a product state, so it is
no longer entangled: the projective measurement has destroyed the entanglement.
(Consistency check: if the outcome is not recorded, the first qubit is in
$\frac58\cdot\frac15\begin{pmatrix}4&2\\2&1\end{pmatrix}+\frac38\begin{pmatrix}0&0\\0&1\end{pmatrix}=\frac12\begin{pmatrix}1&1/2\\1/2&1\end{pmatrix}=\rho_A$,
unchanged, as it must be for a measurement performed on the other qubit.)
\end{enumerate}

\end{document}
